\documentclass[12pt,a4paper]{article}
\usepackage[utf8]{inputenc}
\usepackage[spanish,es-noquoting]{babel}
\usepackage[T1]{fontenc}
\usepackage{lmodern}
\usepackage{geometry}
\geometry{top=2cm, bottom=2cm, left=2.5cm, right=2.5cm}
\usepackage{xcolor}
\usepackage{listings}
\usepackage{enumitem}
\usepackage{fancyhdr}
\usepackage{hyperref}

\definecolor{codegreen}{rgb}{0,0.6,0}
\definecolor{backcolour}{rgb}{0.95,0.95,0.92}

\lstset{
    language=C,
    backgroundcolor=\color{backcolour},
    commentstyle=\color{codegreen},
    keywordstyle=\color{blue},
    basicstyle=\ttfamily\small,
    breaklines=true,
    frame=single,
    numbers=left,
    tabsize=4,
    extendedchars=true,
    showstringspaces=false,
    literate={á}{{\'a}}1 {é}{{\'e}}1 {í}{{\'i}}1 {ó}{{\'o}}1 {ú}{{\'u}}1 {ñ}{{\~n}}1
}

\pagestyle{fancy}
\fancyhf{}
\lhead{Monitorías Sistemas Operativos 2026-I}
\rhead{Capítulo 0: Funciones}
\cfoot{\thepage}

\title{\textbf{Guía de Aprendizaje: Funciones en C}}
\author{Monitorías de Sistemas Operativos}
\date{}

\begin{document}

\maketitle

\section{Motivación}
Construir un Sistema Operativo completo en un solo bloque de código sería imposible. Las funciones nos permiten dividir problemas grandes en partes pequeñas, reutilizables y fáciles de probar.

\section{Idea Fundamental}
Una función es un bloque de código independiente que realiza una tarea específica y puede devolver un resultado.

\section{Sintaxis Básica}
\begin{lstlisting}[numbers=none]
tipo_retorno nombre_funcion(parametros){
    // Cuerpo
    return valor;
}
\end{lstlisting}

\section{Ejemplo Mínimo Funcional}
\begin{lstlisting}
#include <stdio.h>

int cuadrado(int n){
    return n * n;
}

int main(){
    int x = 5;
    printf("El cuadrado es %d\n", cuadrado(x));
    return 0;
}
\end{lstlisting}

\section{Explicación Paso a Paso}
\begin{enumerate}
    \item Se define la función \texttt{cuadrado} que recibe un entero y devuelve un entero.
    \item En \texttt{main}, se llama a la función pasando \texttt{x} como argumento.
    \item La ejecución salta a la función, calcula el resultado y lo devuelve.
    \item El valor devuelto se usa en el \texttt{printf}.
\end{enumerate}

\section{Qué ocurre en memoria}
Cada vez que se llama a una función, se crea un \textbf{Stack Frame} en la memoria Stack. Allí se guardan los parámetros y las variables locales. Al terminar la función, ese espacio se libera automáticamente.

\section{Patrones comunes de uso}
Modularización para cálculos repetitivos:
\begin{lstlisting}
float calcular_promedio(int suma, int cantidad){
    return (float)suma / cantidad;
}
\end{lstlisting}

\section{Errores Comunes}
\begin{itemize}
    \item \textbf{Variables locales inaccesibles:} Intentar usar una variable de una función desde otra.
    \item \textbf{Falta de prototipo:} Llamar a una función antes de definirla sin haber declarado su prototipo arriba.
\end{itemize}

\section{Buenas Prácticas}
\begin{itemize}
    \item Una función debe hacer solo una cosa (\textit{Single Responsibility Principle}).
    \item Usar verbos para los nombres de las funciones (ej: \texttt{ordenar\_arreglo}).
\end{itemize}

\section{Ejercicios}
\begin{enumerate}
    \item Escriba una función que reciba el radio de un círculo y devuelva su área.
    \item Cree una función que reciba un número e imprima si es par o impar (sin devolver valor).
\end{enumerate}

\end{document}
